% Abstract for the 2nd INFN-IN2P3 bi-national conference on detector R&D (LNF Frascati, 2-4 Dec 2026).
% Text kept in sync with abstract.md (Draft v2) and abstract.txt (plain text for email).
\documentclass[11pt,a4paper]{article}
\usepackage[margin=2.5cm]{geometry}
\usepackage[T1]{fontenc}
\usepackage[utf8]{inputenc}
\usepackage{lmodern}
\usepackage{siunitx}
\usepackage{microtype}
\pagestyle{empty}

\begin{document}

\begin{center}
  {\Large\bfseries A resistive-strip Micromegas TPC for the X17 search at n\_TOF:\\[2pt]
  design, track reconstruction and performance\par}
  \vspace{10pt}
  % TODO: co-authors and affiliations
  {\large D.~Neff\par}
  {\small CEA Saclay, Irfu\par}
  \vspace{4pt}
  {\small Session: Gas detectors\par}
\end{center}

\vspace{6pt}

The X17 search at the CERN n\_TOF neutron facility requires tracking $e^+e^-$ pairs from neutron capture on
$^3$He. To test the anomaly, their opening angle must be measured with as little material as possible. We
developed a resistive-strip bulk Micromegas that works as a compact TPC: one \SI{30}{mm} drift gap gives a
three-dimensional track segment. Its $40\times\SI{40}{cm^2}$ area is read out through a pad matrix with a
\SI{0.78}{mm} pitch, capacitively coupled to two orthogonal strip views.

The resistive layer presents unique challenges for TPC operation. Charge spreads to the
neighbouring strips with an RC delay, and this distorts the arrival time measured on each strip: angles taken
from those times read about \ang{4} too steep. Our forward model fits all the waveforms in a cluster at once,
using the measured sharing kernel, and removes the bias.

The chamber reaches an angular resolution of \ang{1.2}, \SI{93}{\percent} efficiency and a spatial resolution
of \SI{180}{\micro m}. At n\_TOF, the reconstructed tracks image the $^3$He target capsule, demonstrating the
detector's efficacy as a TPC tracker.

\end{document}
